\section*{附录 C\quad 关键算法代码}
\begingroup\normalsize\songti
\definecolor{codebg}{RGB}{247,248,250}
\definecolor{codehl}{RGB}{202,232,248}
\lstset{language=Python,basicstyle=\ttfamily\fontsize{12pt}{15.6pt}\selectfont,backgroundcolor=\color{codebg},columns=fullflexible,keepspaces=true,showstringspaces=false,breaklines=true,breakatwhitespace=false,numbers=none,frame=none,xleftmargin=6pt,xrightmargin=6pt,aboveskip=8pt,belowskip=10pt,linebackgroundcolor={\ifcsname CodeHL@\thelstnumber\endcsname\color{codehl}\else\color{codebg}\fi},linebackgroundsep=0pt,linebackgroundwidth=\linewidth}

\par\Needspace{5\baselineskip}
\noindent\textbf{代码 C1\quad 媒体时间支持与公共时间窗构造}\par
\begingroup
\expandafter\def\csname CodeHL@14\endcsname{}
\expandafter\def\csname CodeHL@16\endcsname{}
\expandafter\def\csname CodeHL@17\endcsname{}
\expandafter\def\csname CodeHL@18\endcsname{}
\expandafter\def\csname CodeHL@31\endcsname{}
\expandafter\def\csname CodeHL@33\endcsname{}
\lstinputlisting[firstnumber=1]{appendices/expanded_code/code_01.py}
\endgroup

\par\Needspace{5\baselineskip}
\noindent\textbf{代码 C2\quad 最大时间重叠与原生来源选择}\par
\begingroup
\expandafter\def\csname CodeHL@18\endcsname{}
\expandafter\def\csname CodeHL@20\endcsname{}
\expandafter\def\csname CodeHL@22\endcsname{}
\expandafter\def\csname CodeHL@24\endcsname{}
\expandafter\def\csname CodeHL@25\endcsname{}
\lstinputlisting[firstnumber=1]{appendices/expanded_code/code_02.py}
\endgroup

\par\Needspace{5\baselineskip}
\noindent\textbf{代码 C3\quad 掩码均值、融合与双任务输出}\par
\begingroup
\expandafter\def\csname CodeHL@4\endcsname{}
\expandafter\def\csname CodeHL@5\endcsname{}
\expandafter\def\csname CodeHL@21\endcsname{}
\expandafter\def\csname CodeHL@27\endcsname{}
\expandafter\def\csname CodeHL@28\endcsname{}
\expandafter\def\csname CodeHL@35\endcsname{}
\expandafter\def\csname CodeHL@37\endcsname{}
\expandafter\def\csname CodeHL@38\endcsname{}
\expandafter\def\csname CodeHL@47\endcsname{}
\expandafter\def\csname CodeHL@50\endcsname{}
\expandafter\def\csname CodeHL@51\endcsname{}
\expandafter\def\csname CodeHL@52\endcsname{}
\lstinputlisting[firstnumber=1]{appendices/expanded_code/code_03.py}
\endgroup

\par\Needspace{5\baselineskip}
\noindent\textbf{代码 C4\quad 内容加权残差池化}\par
\begingroup
\expandafter\def\csname CodeHL@8\endcsname{}
\expandafter\def\csname CodeHL@9\endcsname{}
\expandafter\def\csname CodeHL@20\endcsname{}
\expandafter\def\csname CodeHL@34\endcsname{}
\expandafter\def\csname CodeHL@44\endcsname{}
\expandafter\def\csname CodeHL@46\endcsname{}
\lstinputlisting[firstnumber=1]{appendices/expanded_code/code_04.py}
\endgroup

\par\Needspace{5\baselineskip}
\noindent\textbf{代码 C5\quad 连续缺失区间与有效位语义}\par
\begingroup
\expandafter\def\csname CodeHL@4\endcsname{}
\expandafter\def\csname CodeHL@13\endcsname{}
\expandafter\def\csname CodeHL@15\endcsname{}
\expandafter\def\csname CodeHL@19\endcsname{}
\expandafter\def\csname CodeHL@21\endcsname{}
\expandafter\def\csname CodeHL@24\endcsname{}
\expandafter\def\csname CodeHL@29\endcsname{}
\expandafter\def\csname CodeHL@30\endcsname{}
\expandafter\def\csname CodeHL@53\endcsname{}
\expandafter\def\csname CodeHL@54\endcsname{}
\lstinputlisting[firstnumber=1]{appendices/expanded_code/code_05.py}
\endgroup

\par\Needspace{5\baselineskip}
\noindent\textbf{代码 C6\quad 固定缺失场景与统一评价}\par
\begingroup
\expandafter\def\csname CodeHL@7\endcsname{}
\expandafter\def\csname CodeHL@10\endcsname{}
\expandafter\def\csname CodeHL@12\endcsname{}
\expandafter\def\csname CodeHL@15\endcsname{}
\expandafter\def\csname CodeHL@50\endcsname{}
\expandafter\def\csname CodeHL@62\endcsname{}
\expandafter\def\csname CodeHL@65\endcsname{}
\expandafter\def\csname CodeHL@74\endcsname{}
\expandafter\def\csname CodeHL@80\endcsname{}
\lstinputlisting[firstnumber=1]{appendices/expanded_code/code_06.py}
\endgroup

\par\Needspace{5\baselineskip}
\noindent\textbf{代码 C7\quad 双任务损失与参数更新循环}\par
\begingroup
\expandafter\def\csname CodeHL@2\endcsname{}
\expandafter\def\csname CodeHL@9\endcsname{}
\expandafter\def\csname CodeHL@32\endcsname{}
\expandafter\def\csname CodeHL@36\endcsname{}
\expandafter\def\csname CodeHL@40\endcsname{}
\expandafter\def\csname CodeHL@45\endcsname{}
\lstinputlisting[firstnumber=1]{appendices/expanded_code/code_07.py}
\endgroup

\par\Needspace{5\baselineskip}
\noindent\textbf{代码 C8\quad 评价指标与综合得分}\par
\begingroup
\expandafter\def\csname CodeHL@1\endcsname{}
\expandafter\def\csname CodeHL@6\endcsname{}
\expandafter\def\csname CodeHL@12\endcsname{}
\expandafter\def\csname CodeHL@15\endcsname{}
\expandafter\def\csname CodeHL@17\endcsname{}
\lstinputlisting[firstnumber=1]{appendices/expanded_code/code_08.py}
\endgroup

\par\Needspace{5\baselineskip}
\noindent\textbf{代码 C9\quad 固定预测器输入与参数加载}\par
\begingroup
\expandafter\def\csname CodeHL@15\endcsname{}
\expandafter\def\csname CodeHL@36\endcsname{}
\expandafter\def\csname CodeHL@37\endcsname{}
\expandafter\def\csname CodeHL@40\endcsname{}
\expandafter\def\csname CodeHL@41\endcsname{}
\expandafter\def\csname CodeHL@43\endcsname{}
\expandafter\def\csname CodeHL@49\endcsname{}
\expandafter\def\csname CodeHL@50\endcsname{}
\expandafter\def\csname CodeHL@52\endcsname{}
\lstinputlisting[firstnumber=1]{appendices/expanded_code/code_09.py}
\endgroup

\par\Needspace{5\baselineskip}
\noindent\textbf{代码 C10\quad 模态联盟、精确贡献与成对交互}\par
\begingroup
\expandafter\def\csname CodeHL@21\endcsname{}
\expandafter\def\csname CodeHL@46\endcsname{}
\expandafter\def\csname CodeHL@49\endcsname{}
\expandafter\def\csname CodeHL@51\endcsname{}
\expandafter\def\csname CodeHL@64\endcsname{}
\expandafter\def\csname CodeHL@89\endcsname{}
\expandafter\def\csname CodeHL@91\endcsname{}
\expandafter\def\csname CodeHL@93\endcsname{}
\expandafter\def\csname CodeHL@95\endcsname{}
\expandafter\def\csname CodeHL@96\endcsname{}
\expandafter\def\csname CodeHL@98\endcsname{}
\expandafter\def\csname CodeHL@101\endcsname{}
\lstinputlisting[firstnumber=1]{appendices/expanded_code/code_10.py}
\endgroup

\par\Needspace{5\baselineskip}
\noindent\textbf{代码 C11\quad 连续遮挡与关键区间筛选}\par
\begingroup
\expandafter\def\csname CodeHL@8\endcsname{}
\expandafter\def\csname CodeHL@43\endcsname{}
\expandafter\def\csname CodeHL@50\endcsname{}
\expandafter\def\csname CodeHL@58\endcsname{}
\expandafter\def\csname CodeHL@64\endcsname{}
\expandafter\def\csname CodeHL@67\endcsname{}
\expandafter\def\csname CodeHL@70\endcsname{}
\expandafter\def\csname CodeHL@71\endcsname{}
\expandafter\def\csname CodeHL@74\endcsname{}
\expandafter\def\csname CodeHL@85\endcsname{}
\lstinputlisting[firstnumber=1]{appendices/expanded_code/code_11.py}
\endgroup

\par\Needspace{5\baselineskip}
\noindent\textbf{代码 C12\quad 删除曲线与按视频分组区间估计}\par
\begingroup
\expandafter\def\csname CodeHL@4\endcsname{}
\expandafter\def\csname CodeHL@11\endcsname{}
\expandafter\def\csname CodeHL@13\endcsname{}
\expandafter\def\csname CodeHL@15\endcsname{}
\expandafter\def\csname CodeHL@16\endcsname{}
\expandafter\def\csname CodeHL@17\endcsname{}
\expandafter\def\csname CodeHL@22\endcsname{}
\expandafter\def\csname CodeHL@23\endcsname{}
\expandafter\def\csname CodeHL@24\endcsname{}
\expandafter\def\csname CodeHL@25\endcsname{}
\expandafter\def\csname CodeHL@26\endcsname{}
\expandafter\def\csname CodeHL@28\endcsname{}
\expandafter\def\csname CodeHL@31\endcsname{}
\expandafter\def\csname CodeHL@34\endcsname{}
\expandafter\def\csname CodeHL@35\endcsname{}
\expandafter\def\csname CodeHL@38\endcsname{}
\expandafter\def\csname CodeHL@41\endcsname{}
\expandafter\def\csname CodeHL@43\endcsname{}
\expandafter\def\csname CodeHL@47\endcsname{}
\expandafter\def\csname CodeHL@48\endcsname{}
\expandafter\def\csname CodeHL@51\endcsname{}
\expandafter\def\csname CodeHL@52\endcsname{}
\expandafter\def\csname CodeHL@53\endcsname{}
\expandafter\def\csname CodeHL@55\endcsname{}
\expandafter\def\csname CodeHL@56\endcsname{}
\expandafter\def\csname CodeHL@57\endcsname{}
\expandafter\def\csname CodeHL@58\endcsname{}
\expandafter\def\csname CodeHL@59\endcsname{}
\expandafter\def\csname CodeHL@60\endcsname{}
\expandafter\def\csname CodeHL@61\endcsname{}
\expandafter\def\csname CodeHL@62\endcsname{}
\expandafter\def\csname CodeHL@63\endcsname{}
\expandafter\def\csname CodeHL@64\endcsname{}
\expandafter\def\csname CodeHL@73\endcsname{}
\lstinputlisting[firstnumber=1]{appendices/expanded_code/code_12.py}
\endgroup

\par\Needspace{5\baselineskip}
\noindent\textbf{代码 C13\quad 文本词元与原文字符来源映射}\par
\begingroup
\expandafter\def\csname CodeHL@27\endcsname{}
\expandafter\def\csname CodeHL@33\endcsname{}
\expandafter\def\csname CodeHL@35\endcsname{}
\expandafter\def\csname CodeHL@40\endcsname{}
\expandafter\def\csname CodeHL@44\endcsname{}
\expandafter\def\csname CodeHL@46\endcsname{}
\expandafter\def\csname CodeHL@49\endcsname{}
\expandafter\def\csname CodeHL@50\endcsname{}
\expandafter\def\csname CodeHL@51\endcsname{}
\expandafter\def\csname CodeHL@52\endcsname{}
\lstinputlisting[firstnumber=1]{appendices/expanded_code/code_13.py}
\endgroup
\endgroup
